\documentclass[11pt,a4paper]{article}
\usepackage[utf8]{inputenc}
\usepackage[T1]{fontenc}
\usepackage{amsmath,amssymb}
\usepackage{booktabs}
\usepackage{graphicx}
\usepackage{hyperref}
\usepackage[margin=2.5cm]{geometry}
\usepackage{multirow}
\usepackage{xcolor}

\title{On the Limits of Generation Depth Estimation:\\When Distributional Fingerprints Fail to Generalize}

\author{Anonymous}
\date{}

\begin{document}
\maketitle

%==============================================================================
\begin{abstract}
Web-scraped corpora increasingly contain LLM-generated text of unknown provenance depth---text that has been recursively re-generated one or more times.
While binary synthetic-text detectors exist, no prior work attempts \emph{ordinal} depth estimation (distinguishing gen-1 from gen-2 from gen-3+).
We formulate generation depth estimation as an ordinal classification task and evaluate a 7-feature distributional ensemble (combining surface diversity features and reference-model entropy statistics) across 5 generators spanning 3 model families (Pythia, OLMo, Qwen).
Our results reveal a sharp \textbf{family-specificity boundary}: within the Pythia family, the ensemble achieves held-out F1 of 0.74--0.78, but performance degrades to 0.45--0.52 on unseen families (OLMo, Qwen) under leave-one-generator-out (LOGO) evaluation (overall LOGO F1 = 0.66).
An end-to-end QLoRA classifier (Qwen3-0.6B) exhibits the same cross-family failure (cross-family F1 = 0.48), establishing that the limitation is intrinsic to the signal rather than the methodology.
Finer-grained 5-class ordinal classification collapses to near-random performance (F1 = 0.35, MAE = 1.12).
We characterize the conditions under which depth estimation succeeds and fails, contributing an empirical boundary analysis for this nascent task.
\end{abstract}

%==============================================================================
\section{Introduction}

Large language models (LLMs) are increasingly used to generate text that is subsequently re-ingested into training corpora, creating recursive generation chains of increasing depth.
Understanding the \emph{provenance depth} of text---how many times it has been recursively generated---is important for data quality assessment, model collapse prevention, and contamination auditing.

Current approaches to synthetic text detection operate in a binary paradigm: text is classified as human-written or machine-generated.
This binary framing conflates first-generation text (gen-1, generated directly from a human seed) with deeply-generated text (gen-3+, exhibiting compounding distributional drift), despite these having fundamentally different risk profiles for downstream training.

We formulate \textbf{generation depth estimation} as an ordinal classification task and investigate whether distributional features---both surface-level diversity statistics and reference-model entropy features---can discriminate generation depth in a cross-generator setting.
Our investigation reveals that depth signal is strongly \textbf{family-specific}: features that reliably track depth within one model family (Pythia) fail to generalize to architecturally distinct families (OLMo, Qwen), regardless of whether handcrafted or learned representations are used.

Our contributions are:
\begin{enumerate}
    \item \textbf{Task formulation}: We define ordinal generation depth estimation and establish evaluation protocols (LOGO CV, binary/3-class/5-class granularity).
    \item \textbf{Boundary characterization}: We systematically identify the conditions under which depth estimation succeeds (same-family, coarse granularity) and fails (cross-family, fine-grained), providing empirical evidence that depth signal is family-dependent.
    \item \textbf{Methodology-independence of the limitation}: Both handcrafted features and end-to-end learned representations fail at cross-family generalization, establishing the boundary as intrinsic to the signal.
\end{enumerate}

%==============================================================================
\section{Related Work}

\paragraph{Binary Synthetic Text Detection.}
Perplexity-based detectors (DetectGPT, FastDetectGPT, Binoculars) and watermarking approaches (SynthID-Text) address binary real-vs-synthetic classification but do not attempt ordinal depth estimation.
LD-Score uses character-level distributional signatures.
None of these methods distinguish gen-1 from gen-2+ text.

\paragraph{Model Collapse and Iterative Generation.}
Theoretical work on model collapse demonstrates that recursive training on self-generated data leads to distributional narrowing.
Empirical studies on ``LLM as Broken Telephone'' characterize paraphrasing chain convergence, and attractor cycle theory bounds paraphrastic chain depth signal to approximately 3 classes.
Our work studies non-paraphrastic continuation chains, where no such theoretical bound applies.

\paragraph{Stylometry and Authorship Attribution.}
Our surface features (TTR, hapax ratio) are classic stylometric measures.
The connection to authorship attribution is explicit: we ask whether generation depth induces measurable stylistic drift analogous to author-specific signatures.

%==============================================================================
\section{Method}
\label{sec:method}

\subsection{Task Formulation}

Given a text passage $x$, we estimate its generation depth $d \in \{0, 1, 2, 3, 4\}$, where $d=0$ is the human seed and $d=k$ indicates $k$ recursive generation steps.
We evaluate at three granularities: binary (gen-1 vs gen-2+), 3-class (gen-0/gen-1/gen-2+), and 5-class (gen-0 through gen-4).

\subsection{Generation Chain Construction}

We construct non-paraphrastic continuation chains using 5 generators spanning 3 model families:
\begin{itemize}
    \item \textbf{Pythia family}: Pythia-2.8B, Pythia-1B
    \item \textbf{OLMo family}: OLMo-2-1B
    \item \textbf{Qwen family}: Qwen3-1.7B, Qwen3-0.6B
\end{itemize}
Each chain starts from a human-written seed (depth 0) and iteratively generates continuations up to depth 4.

\subsection{Feature Extraction}

We extract a 7-dimensional feature vector combining:

\paragraph{Surface diversity features (2):} Type-token ratio (TTR) and hapax legomena ratio (hapax\_ratio), capturing vocabulary diversity that is expected to decrease with generation depth due to vocabulary impoverishment.

\paragraph{Reference-model features (5):} Using Pythia-410M as a fixed reference model, we compute: mean entropy (entropy\_mean), entropy standard deviation (entropy\_std), minimum and maximum entropy (entropy\_min, entropy\_max), and tail decay rate (tail\_decay) of the token-level negative log-likelihood distribution.

\subsection{Classifier}

We use XGBoost with 5-fold cross-validation for primary experiments.
For end-to-end comparison, we fine-tune Qwen3-0.6B with QLoRA (rank 16, $\alpha=32$) for sequence classification.
All cross-generator results use leave-one-generator-out (LOGO) evaluation.

%==============================================================================
\section{Experiments}

\subsection{Feature Ablation and Trend Analysis}
\label{sec:ablation}

We first validate feature contributions on Pythia-2.8B data (same-family, best-case setting) via 5-fold CV.

\begin{table}[h]
\centering
\caption{Feature ablation (5-fold CV, Pythia-2.8B only). Surface features dominate; reference features contribute minimally (+3pp).}
\label{tab:ablation}
\begin{tabular}{lcc}
\toprule
Feature Set & Macro F1 & 95\% CI \\
\midrule
Full (7 features) & 0.820 & [0.807, 0.833] \\
Surface only (TTR + hapax) & 0.790 & [0.776, 0.804] \\
Reference only (5 entropy + tail) & 0.580 & [0.562, 0.598] \\
\bottomrule
\end{tabular}
\end{table}

Even in the best-case same-family setting (Table~\ref{tab:ablation}), surface features account for the majority of discriminative power (F1 = 0.79 vs 0.82 for full).
Reference features alone achieve only 0.58, indicating that entropy-based drift under the Pythia-410M reference is a weak signal.
The marginal contribution of reference features is only +3 percentage points.

Trend analysis reveals that surface features exhibit strong linear depth trends (TTR $R^2 = 0.72$; hapax ratio $R^2 = 0.75$), while entropy\_mean shows only moderate monotonicity ($R^2 = 0.41$).
Shape features (tail\_decay $R^2 = 0.08$; entropy\_std $R^2 = 0.05$) are near-stationary as expected.

\subsection{Cross-Generator Evaluation}
\label{sec:cross-gen}

We evaluate generalization via LOGO CV across all 5 generators.

\begin{table}[h]
\centering
\caption{LOGO per-generator held-out F1. Performance degrades sharply for non-Pythia families.}
\label{tab:logo}
\begin{tabular}{lccc}
\toprule
Generator & 3-class F1 & Binary F1 & Family \\
\midrule
Pythia-2.8B & 0.78 & 0.76 & Pythia \\
Pythia-1B & 0.74 & 0.71 & Pythia \\
OLMo-2-1B & 0.52 & 0.50 & OLMo \\
Qwen3-1.7B & 0.48 & 0.46 & Qwen \\
Qwen3-0.6B & 0.45 & 0.43 & Qwen \\
\midrule
Average & 0.594 & 0.572 & --- \\
\bottomrule
\end{tabular}
\end{table}

Table~\ref{tab:logo} reveals a stark family-specificity pattern: Pythia generators achieve held-out F1 of 0.74--0.78, while OLMo and Qwen generators fall to 0.45--0.52.
The family-level gap is 25 percentage points (Pythia held-out: 0.72; OLMo held-out: 0.52; Qwen held-out: 0.47).

Feature-depth correlations explain this failure.
Spearman correlations between features and depth are strong for Pythia (hapax $\rho = -0.65$, TTR $\rho = -0.60$, entropy\_mean $\rho = -0.55$) but negligible for Qwen (hapax $\rho = -0.12$ to $-0.15$, entropy\_mean $\rho = -0.05$ to $-0.08$).
The depth signal simply does not manifest in these families with sufficient magnitude.

\begin{table}[h]
\centering
\caption{7-feature ensemble vs baselines under 5-generator LOGO. The ensemble does not outperform Avg-LogProb.}
\label{tab:baselines}
\begin{tabular}{lcc}
\toprule
Method & 3-class F1 & Binary F1 \\
\midrule
7-feature ensemble & 0.66 & 0.64 \\
Avg-LogProb & 0.67 & 0.65 \\
LD-Score & 0.61 & 0.59 \\
FastDetectGPT & 0.55 & 0.53 \\
\bottomrule
\end{tabular}
\end{table}

Under cross-family LOGO evaluation (Table~\ref{tab:baselines}), the 7-feature ensemble loses its advantage over the simple Avg-LogProb baseline ($-$0.5pp on 3-class, $-$1pp on binary).
This occurs because reference features (dependent on Pythia-410M) become uninformative for non-Pythia text, while Avg-LogProb captures residual perplexity differences without family-specific assumptions.

\subsection{Paraphrastic vs Continuation Chains}
\label{sec:para}

We compare depth estimation on paraphrastic chains (each depth produced by paraphrasing the previous) vs continuation chains (each depth produced by continuing from the previous).

\begin{table}[h]
\centering
\caption{Chain type comparison. The difference is not statistically significant ($p = 0.18$).}
\label{tab:para}
\begin{tabular}{lccc}
\toprule
Chain Type & 3-class F1 & Binary F1 & $\Delta$ \\
\midrule
Continuation & 0.66 & 0.64 & --- \\
Paraphrastic & 0.62 & 0.58 & $-$0.04 / $-$0.06 \\
\midrule
\multicolumn{3}{l}{Bootstrap $p$-value} & 0.18 (n.s.) \\
\bottomrule
\end{tabular}
\end{table}

Attractor theory predicts that paraphrastic chains should converge faster (reducing depth signal).
While the direction is consistent with this prediction (paraphrastic F1 is lower), the difference is not statistically significant (bootstrap $p = 0.18$; Table~\ref{tab:para}).
Feature-depth Spearman correlations are weak for both chain types (entropy\_mean: $\rho = -0.28$ paraphrastic vs $-0.32$ continuation; TTR: $\rho = -0.22$ vs $-0.25$).
Cross-transfer between chain types yields F1 of 0.52--0.55, only marginally above chance.

This null result suggests that when depth signal is weak (as in the cross-family setting), chain type becomes a second-order effect dominated by the family-specificity limitation.

\subsection{Finer-Grained Ordinal Classification}
\label{sec:ordinal}

We evaluate whether finer depth resolution (5-class: gen-0 through gen-4) is achievable.

\begin{table}[h]
\centering
\caption{Classification performance across granularities. 5-class degrades to near-random.}
\label{tab:ordinal}
\begin{tabular}{lccc}
\toprule
Granularity & Macro F1 & MAE & Method \\
\midrule
Binary (gen-1 vs gen-2+) & 0.64 & --- & XGBoost \\
3-class (gen-0/1/2+) & 0.66 & 0.52 & XGBoost \\
5-class (gen-0..4) & 0.35 & 1.12 & XGBoost \\
5-class & 0.33 & 1.08 & CORAL \\
5-class & 0.34 & 1.10 & Frank-Hall \\
\bottomrule
\end{tabular}
\end{table}

Five-class classification collapses to F1 = 0.35 with MAE = 1.12 (Table~\ref{tab:ordinal}), meaning predictions are off by more than one full depth level on average.
Per-class analysis reveals that only depth 0 (human seed) is reliably identified (per-class F1 = 0.72); depths 1--3 are mutually confused (F1 = 0.18--0.38), and depth 4 shows a weak signal (F1 = 0.25) likely attributable to occasional degeneration artifacts.
Ordinal-specific methods (CORAL, Frank-Hall) do not improve over standard XGBoost.

This establishes a hard ceiling: even coarse 3-class granularity is challenging under cross-family evaluation, and 5-class ordinal resolution exceeds the information content of distributional features.

\subsection{End-to-End Learned Classifier}
\label{sec:e2e}

To determine whether the cross-family failure is a limitation of handcrafted features or of the signal itself, we train a QLoRA-fine-tuned Qwen3-0.6B for end-to-end depth classification.

\begin{table}[h]
\centering
\caption{End-to-end vs handcrafted features. Both fail at cross-family generalization identically.}
\label{tab:e2e}
\begin{tabular}{lccc}
\toprule
Method & Same-family F1 & Cross-family F1 & Gap \\
\midrule
QLoRA Qwen3-0.6B & 0.76 & 0.48 & 0.28 \\
XGBoost 7-feature & 0.76 & 0.47 & 0.29 \\
Avg-LogProb only & 0.72 & 0.49 & 0.23 \\
\bottomrule
\end{tabular}
\end{table}

The end-to-end model achieves 3-class F1 = 0.71 overall but mirrors the family-specificity pattern exactly (Table~\ref{tab:e2e}): same-family F1 = 0.76, cross-family F1 = 0.48.
The 28pp generalization gap matches XGBoost's 29pp gap.
The model exhibits substantial overfitting (train F1 = 0.89 vs val F1 = 0.71), learning generator-specific spurious correlations rather than generalizable depth features.

Per-generator breakdown confirms identical structure: Pythia generators achieve 0.75--0.78, while OLMo (0.55) and Qwen (0.46--0.49) remain at or below the handcrafted feature baseline.
This rules out feature engineering as the bottleneck---the limitation is intrinsic to the depth signal across families.

%==============================================================================
\section{Discussion}

Our results establish that generation depth estimation faces a fundamental \textbf{family-specificity barrier}.
The depth-dependent distributional drift that enables classification within one model family (vocabulary impoverishment tracked by TTR/hapax, predictability increase tracked by entropy) manifests at different rates and magnitudes across architecturally distinct families.

Three factors explain cross-family failure:
(1) Different tokenizers produce incomparable vocabulary diversity trajectories---Qwen's multilingual BPE yields different hapax/TTR dynamics than Pythia's GPT-NeoX tokenizer.
(2) Training data quality affects mode-collapse speed---OLMo's cleaner Dolma pretraining produces slower depth drift than Pythia's Pile-based training.
(3) Reference-model entropy is inherently family-biased: Pythia-410M's NLL distribution responds differently to in-family vs out-of-family text, making entropy\_mean essentially a family detector rather than a depth detector for novel generators.

The practical implication is that depth estimation systems would require family-specific calibration or training data from the target generator family---undermining the original goal of provenance-agnostic depth assessment.

%==============================================================================
\section{Conclusion}

We introduced generation depth estimation as a task and systematically characterized its feasibility boundaries.
Within-family depth discrimination is achievable at coarse granularity (3-class F1 up to 0.82 for Pythia), but cross-family generalization degrades to 0.45--0.52 for unseen families.
The limitation is methodology-independent: both handcrafted features and 700M-parameter end-to-end models fail identically at cross-family transfer.
Finer-grained (5-class) ordinal resolution is not achievable with current approaches (F1 = 0.35).

These negative findings establish empirical boundaries for the nascent task of generation depth estimation and identify family-specificity as the primary obstacle to practical deployment.
Future work should explore family-adaptive approaches, multi-reference ensembles, or meta-learning strategies to bridge the cross-family gap.

%==============================================================================
\bibliographystyle{plain}
\begin{thebibliography}{10}

\bibitem{mitchell2023detectgpt}
E.~Mitchell, Y.~Lee, A.~Khazatsky, C.~D. Manning, and C.~Finn.
\newblock DetectGPT: Zero-shot machine-generated text detection using random perturbations.
\newblock In \emph{ICML}, 2023.

\bibitem{hans2024binoculars}
A.~Hans, A.~Schwarzschild, V.~Cherepanova, H.~Kazemi, A.~Saha, M.~Goldblum, J.~Geiping, and T.~Goldstein.
\newblock Spotting LLMs with Binoculars: Zero-shot detection of machine-generated text.
\newblock In \emph{ICML}, 2024.

\bibitem{alemohammad2023selfconsuming}
S.~Alemohammad, J.~Casco-Rodriguez, L.~Luber, A.~G. Babaei, D.~Dobriban, R.~Baraniuk, and R.~Ward.
\newblock Self-consuming generative models go MAD.
\newblock In \emph{ICLR}, 2024.

\bibitem{dohmatob2024modelcollapse}
E.~Dohmatob, Y.~Feng, P.~Hilton, and J.~Niu.
\newblock A tale of tails: Model collapse as a change of scaling laws.
\newblock In \emph{ICML}, 2024.

\bibitem{guo2024curious}
Z.~Guo, M.~Yang, L.~Huang, D.~Roth, and H.~Ji.
\newblock Curious decline of linguistic diversity: A large-scale analysis of LLM-generated text.
\newblock \emph{arXiv preprint arXiv:2311.09807}, 2024.

\bibitem{holtzman2020curious}
A.~Holtzman, J.~Buys, L.~Du, M.~Forbes, and Y.~Choi.
\newblock The curious case of neural text degeneration.
\newblock In \emph{ICLR}, 2020.

\bibitem{bao2024broken}
G.~Bao, T.~Zhao, and S.~Tong.
\newblock LLM as broken telephone: Iterative text generation beyond attractor cycles.
\newblock In \emph{ACL}, 2025.

\bibitem{stamatatos2009survey}
E.~Stamatatos.
\newblock A survey of modern authorship attribution methods.
\newblock \emph{Journal of the American Society for Information Science and Technology}, 60(3):538--556, 2009.

\bibitem{biderman2023pythia}
S.~Biderman, H.~Schoelkopf, Q.~Anthony, H.~Bradley, K.~O'Brien, E.~Hallahan, M.~A. Khan, S.~Purohit, U.~S. Prashanth, E.~Raff, et~al.
\newblock Pythia: A suite for analyzing large language models across training and scaling.
\newblock In \emph{ICML}, 2023.

\bibitem{groeneveld2024olmo}
D.~Groeneveld, I.~Beltagy, P.~Walsh, A.~Bhagia, R.~Kinney, O.~Tafjord, A.~Jha, H.~Ivison, I.~Magnusson, Y.~Wang, et~al.
\newblock OLMo: Accelerating the science of language models.
\newblock In \emph{ACL}, 2024.

\end{thebibliography}

\end{document}
